% ============================================================================
% Fichier  : partie4/P04_C13_programme_rapport.tex
% Titre    : Programme de Conformité et Rapport d'Audit
% Auteur   : MINKA MI NGUIDJOÏ Thierry Emmanuel
% Version  : 1.0 -- novembre 2026
% Corrigés : corriges/Corr_C13.tex
% ============================================================================
% Sources : loi n° 2024/017, art. 13, 27(2), 28, 73 (texte officiel) ;
% ISO 19011:2018 (rapport, paraphrasé) ; NRAC C07 (observation en sept
% éléments) ; auditologie.
% Faits nouveaux (bible) : rapport de KORO AUDIT remis le 25/09/2026 ;
% réponse de la direction du 18/09/2026 ; synthèse des six domaines.
% ============================================================================

\chapter{Programme de Conformité et Rapport d'Audit}
\label{chap:programme}
\chaptermark{Programme et rapport}

\epigraph{%
    « Un rapport qui énumère des écarts sans les classer\\
    laisse la direction choisir les plus faciles. »%
}{Principe du cours}

\niveauChapitre{Intermédiaire}{%
    Chapitres~\ref{chap:audit} à~\ref{chap:violations}
}

\begin{boxobjectifs}
\begin{itemize}
    \item Rédiger une observation d'audit en sept éléments.
    \item Synthétiser les domaines d'un audit avec $\mu$, $U$ et $H$ et
          conclure selon la règle de décision.
    \item Prioriser les recommandations selon la gravité pour les
          personnes.
    \item Construire un programme de mise en conformité daté, avec des
          indicateurs.
    \item Auditer la conservation et la purge des données (articles~13
          et~28).
\end{itemize}
\end{boxobjectifs}

\bigskip

% ============================================================================
\section*{Avant de commencer : quarante-six exigences, une page}
% ============================================================================

Yaoundé, lundi 14~septembre~2026. Les papiers de travail de la mission
NYANGA remplissent deux classeurs~: quarante-six exigences cotées dans six
domaines, des dizaines de pièces. Le conseil d'administration, lui, attend
une page. M.~TCHOUA a prévenu~: « Dites-nous ce qui est grave, ce qui coûte
cher et ce que nous devons faire d'abord. »

Mme~BELLO sait qu'une page ne peut pas contenir quarante-six constats. Elle
sait aussi qu'une page qui ne dit pas ce qui reste inconnu tromperait le
conseil. Elle choisit de construire le rapport en trois couches~: une
synthèse d'une page, un constat par domaine avec ses cotes, et les
observations détaillées en annexe.

\begin{boxalert}
\textbf{Violation typique.} Remettre un rapport qui liste les écarts par
ordre d'apparition, sans les relier à des articles, sans dire ce qui reste
inconnu, et sans ordre de priorité.

\smallskip
\textbf{Règle~:} chaque constat se rattache à un article et à une preuve~;
la synthèse dit, par domaine, ce qui est établi et ce qui ne l'est pas~;
les recommandations sont classées par gravité pour les personnes, puis par
exposition de l'organisation.
\end{boxalert}

% ============================================================================
\section{L'Observation en Sept Éléments}
\label{sec:c13-observation}
% ============================================================================

\begin{tableaucours}{cL{2.8cm}L{8.4cm}}{%
    \tableauHeader{N°} &
    \tableauHeader{Élément} &
    \tableauHeader{Contenu}
}
1 & Norme & L'article de la loi, ou l'exigence de la norme, qui fonde le
critère \\
2 & Constatation & Le fait établi, chiffré, avec sa preuve et sa cote \\
3 & Risques & Le risque pour les personnes, puis pour l'organisation \\
4 & Contradictoire & La réponse de l'audité, datée \\
5 & Analyse & La cause de l'écart, avec un niveau de confiance \\
6 & Conclusion & La position de l'auditeur \\
7 & Recommandation & Ce qu'il faut faire, en traitant la cause \\
\end{tableaucours}
\vspace{-0.8em}
\begin{center}{\small\itshape Tableau~13.1. L'observation en sept éléments}\end{center}

\begin{boxscenario}{Mission NYANGA. Observation D2-06~: les oppositions non exécutées}
\textbf{Norme.} Article~40 (droit d'opposition à la prospection)~;
article~64 (délit de traitement maintenu malgré l'opposition).

\textbf{Constatation.} Cinq oppositions à la prospection reçues entre le
23~juin et le 31~août~2026, toutes acceptées par lettre. Trois des cinq
numéros figurent dans la liste d'appels de prospection du 28~août (pièce
D2-06-03). Cote \NON{}.

\textbf{Risques.} Pour les personnes~: sollicitations non voulues, perte de
confiance dans l'exercice de leurs droits. Pour l'établissement~: délit de
l'article~64.

\textbf{Contradictoire} (18~septembre). « Les agents n'ont pas reçu la mise
à jour de la liste. Nous corrigeons. »

\textbf{Analyse} (confiance élevée). La cause n'est pas l'agent~: la liste
d'appels est extraite chaque mois de la base des clients, et aucune
procédure ne reporte les oppositions dans la base.

\textbf{Conclusion.} L'exigence n'est pas satisfaite~; l'écart est
systémique.

\textbf{Recommandation.} Ajouter dans la base un indicateur d'opposition
qui exclut automatiquement le client des extractions de prospection~;
vérifier chaque mois, par sondage, l'absence des opposants dans les listes.
\end{boxscenario}

% ============================================================================
\section{La Synthèse par Domaine}
\label{sec:c13-synthese}
% ============================================================================

\begin{boxlecture}{NYANGA CRÉDIT, synthèse des six domaines, 25 septembre 2026}
\begin{tableaucours}{cL{3.8cm}cccc}{%
    \tableauHeader{Code} &
    \tableauHeader{Domaine} &
    \tableauHeader{O / N / K} &
    \tableauHeader{$\mu$} &
    \tableauHeader{$U$} &
    \tableauHeader{$H$}
}
D1 & Recensement et fondements & 1 / 5 / 2 & 0,13 & 0,25 & 1,30 \\
D2 & Droits des personnes & 1 / 7 / 2 & 0,10 & 0,20 & 1,16 \\
D3 & Sécurité et confidentialité & 4 / 3 / 1 & 0,50 & 0,13 & 1,41 \\
D4 & Sous-traitants et transferts & 1 / 5 / 2 & 0,13 & 0,25 & 1,30 \\
D5 & Violations & 1 / 5 / 2 & 0,13 & 0,25 & 1,30 \\
D6 & Gouvernance & 0 / 3 / 1 & 0 & 0,25 & 0,81 \\
Total & & 8 / 28 / 10 & 0,17 & 0,22 & 1,35 \\
\end{tableaucours}

\medskip
\textbf{Lecture.} Cinq domaines sur six sont opaques au sens de la règle
($U \geq 0{,}20$)~; aucun ne peut être déclaré conforme. Le seul domaine
clair, D3, ne l'est pas non plus~: $\mu = 0{,}50$ atteint le seuil, mais
trois \NON{} portent sur la limitation des accès, et la règle exige
l'absence d'écart grave. Dans tous les domaines, les \NON{} dominent~:
l'opacité n'est pas ce qui empêche de conclure. L'établissement n'est pas
conforme, et l'audit sait pourquoi sur vingt-huit points précis.
\end{boxlecture}

La synthèse distingue trois messages qu'un conseil doit entendre
séparément~: ce qui est établi comme non conforme (vingt-huit écarts), ce
qui reste inconnu et pourquoi (dix \NEK{}, dont six de critère, liés aux
textes d'application et à l'Autorité), et ce qui fonctionne (huit
exigences, dont la moitié en sécurité). Fondre ces trois messages en une
note globale effacerait l'information utile.

% ============================================================================
\section{Prioriser les Recommandations}
\label{sec:c13-priorites}
% ============================================================================

Les recommandations se classent d'abord par \textbf{gravité pour les
personnes}, ensuite par \textbf{exposition de l'organisation} (sanctions
encourues, contentieux), enfin par \textbf{effort}. L'ordre inverse, du
plus facile au plus difficile, est tentant et dangereux~: il traite d'abord
ce qui protège le moins.

\begin{tableaucours}{cL{6.0cm}L{2.8cm}L{2.4cm}}{%
    \tableauHeader{Rang} &
    \tableauHeader{Recommandation} &
    \tableauHeader{Articles} &
    \tableauHeader{Échéance}
}
1 & Informer les 18\,412 personnes de la violation de 2025 & 22 & 30 jours \\
2 & Supprimer le classeur des mauvais payeurs & 9, 36, 48, 66 & Immédiat \\
3 & Exécuter les oppositions et répondre aux demandes en attente & 39, 40,
56, 64 & 30 jours \\
4 & Désactiver le kit de mesure d'audience sans consentement & loi
2010/012, art.~47 & Immédiat \\
5 & Réécrire formulaires et notices, recueillir les consentements & 9, 21 &
90 jours \\
6 & Avenants aux contrats des sous-traitants, droit d'audit & 16, 30 & 90
jours \\
7 & Fixer et appliquer les durées de conservation, purger & 13, 28 & 6
mois \\
8 & Limiter les accès des agences à leurs clients & 27 & 6 mois \\
9 & Documenter les transferts, préparer les autorisations & 19, 32 & 6
mois \\
10 & Désigner un délégué à la protection des données (bonne pratique) &
aucun & 90 jours \\
\end{tableaucours}
\vspace{-0.8em}
\begin{center}{\small\itshape Tableau~13.2. Les dix premières recommandations à NYANGA CRÉDIT}\end{center}

La dixième recommandation n'a pas de fondement légal~; le rapport le dit.
Elle répond à une cause commune à presque tous les écarts~: personne,
dans l'établissement, n'est chargé de la protection des données.

% ============================================================================
\section{La Conservation et la Purge}
\label{sec:c13-conservation}
% ============================================================================

L'article~13 limite la conservation à la durée nécessaire, et l'article~28
renvoie la fixation du délai maximum au référentiel de l'Autorité, « compte
tenu de la finalité du traitement ou de la nature des données collectées ».
En attendant le référentiel, l'organisation fixe elle-même ses durées dans
une politique de conservation, traitement par traitement, en distinguant
trois périodes~: la \textbf{base active} (le temps de la finalité), les
\textbf{archives intermédiaires} (le temps d'une obligation légale ou d'un
contentieux possible, avec accès restreint), et la \textbf{purge}
(effacement, ou anonymisation prouvée pour les statistiques).

L'audit de la conservation compare, pour chaque traitement, la date des
données les plus anciennes à la durée fixée. Le laboratoire APDC-LAB-09
exerce ce calcul sur un extrait de la base des comptes de NYANGA.

\begin{boxscenario}{Mission NYANGA. La politique de conservation proposée}
\begin{tableaucours}{L{3.6cm}L{2.8cm}L{3.2cm}L{2.6cm}}{%
    \tableauHeader{Traitement} &
    \tableauHeader{Base active} &
    \tableauHeader{Archive} &
    \tableauHeader{Justification}
}
Comptes clients & Durée du compte & 10 ans après la clôture & Obligations
comptables (à vérifier) \\
Prospects & 1 an sans contact & Aucune & Finalité seule \\
Candidats non retenus & 6 mois & Aucune & Finalité seule \\
Images des agences & 30 jours & Durée d'une procédure si incident &
Sécurité \\
Appels du service client & 6 mois & Durée d'un litige & Preuve (loi
2010/012, art.~45) \\
\end{tableaucours}

\medskip
\textbf{Constatation} (31~août). 22\,950 comptes clos avant 2016 dépassent
déjà la durée d'archive proposée~; la base des prospects contient des
entrées de 2017.
\end{boxscenario}

% ============================================================================
\section{Le Programme de Mise en Conformité}
\label{sec:c13-programme}
% ============================================================================

Un rapport d'audit se clôt par des recommandations~; un programme les
transforme en projet. Il fixe pour chaque recommandation un responsable,
une échéance, une preuve d'achèvement et un indicateur, puis il prévoit un
audit de suivi.

\begin{tableaucours}{L{5.6cm}L{6.4cm}}{%
    \tableauHeader{Indicateur} &
    \tableauHeader{Définition}
}
Couverture du registre & Traitements découverts et inscrits, sur
traitements existants \\
Traitements documentés pour autorisation & Dossiers complets, sur
traitements du registre \\
Délai médian de réponse aux demandes & Médiane du délai entre réception et
réponse complète \\
Exécution des oppositions & Opposants absents des listes de prospection,
par sondage mensuel \\
Contrats conformes & Contrats aux douze clauses, sur contrats de
sous-traitance \\
Violations informées dans les délais & Violations informées dans le délai
cible, sur violations du registre \\
Données au-delà de la durée & Enregistrements dépassant la durée fixée \\
\end{tableaucours}
\vspace{-0.8em}
\begin{center}{\small\itshape Tableau~13.3. Indicateurs d'un programme de conformité}\end{center}

L'article~27, alinéa~2, impose en outre un rapport annuel à l'Autorité sur
l'état des mesures de sécurité. Le programme peut en préparer la matière
dès maintenant~: les indicateurs de sécurité du programme en sont le
cœur, et le format sera fixé quand le référentiel sera publié.

% ============================================================================
\section{Ce que le Rapport demande à chaque Métier}
\label{sec:c13-metiers}
% ============================================================================

\begin{boxmetier}{Auditeur de la protection des données}
Rédiger chaque observation en sept éléments, synthétiser par domaine avec
les trois mesures, séparer établi, inconnu et satisfaisant, et classer les
recommandations par gravité pour les personnes.
\end{boxmetier}

\begin{boxmetier}{Responsable de traitement}
Répondre au contradictoire par écrit, accepter ou motiver le refus de
chaque recommandation, et nommer un responsable par action.
\end{boxmetier}

\begin{boxmetier}{Délégué à la protection des données}
Piloter le programme, mesurer les indicateurs chaque mois, et préparer
l'audit de suivi.
\end{boxmetier}

\begin{boxresume}
\begin{itemize}
    \item L'observation d'audit compte sept éléments, dont le
          contradictoire et une recommandation qui traite la cause.
    \item La synthèse par domaine sépare ce qui est établi non conforme, ce
          qui est inconnu et pourquoi, et ce qui fonctionne.
    \item Les recommandations se classent par gravité pour les personnes,
          puis par exposition, puis par effort.
    \item La conservation se gère en trois périodes~: base active, archive,
          purge~; l'audit compare l'âge des données aux durées fixées.
    \item À NYANGA, vingt-huit écarts établis, dix \NEK{} dont six de
          critère, et l'information des 18\,412 personnes en première
          recommandation.
\end{itemize}
\end{boxresume}

\begin{boxinfo}
\textbf{Points de contrôle.} Vous avez compris ce chapitre si vous savez~:
\begin{itemize}
    \item rédiger une observation en sept éléments~;
    \item conclure sur un domaine à partir de ses cotes~;
    \item ordonner des recommandations et justifier l'ordre.
\end{itemize}
\end{boxinfo}

\begin{boxexo}[title={\ding{45}\quad Exercice~13.1. Ordonner \niveaubase}]
Classez par gravité pour les personnes~: (a) le registre incomplet~; (b) les
selfies traités dans un pays inconnu~; (c) un contrat de papeterie sans
clause de protection des données~; (d) le pointage biométrique des
salariés~; (e) l'absence de délégué.
\end{boxexo}

\begin{boxexo}[title={\ding{45}\quad Exercice~13.2. Une observation \niveaumoyen}]
Rédigez l'observation D1-04 « le classeur des mauvais payeurs » en sept
éléments, à partir des Chapitres~\ref{chap:cartographie} et~\ref{chap:risque}.
Inventez de façon cohérente la réponse de la direction.
\end{boxexo}

\begin{boxexo}[title={\ding{45}\quad Exercice~13.3. La page du conseil \niveaumoyen}]
Rédigez la synthèse d'une page destinée au conseil d'administration de
NYANGA, sans jargon, en trois paragraphes~: ce qui est établi, ce qui est
inconnu, ce qu'il faut faire d'abord.
\end{boxexo}

\begin{boxexo}[title={\ding{45}\quad Exercice~13.4. L'audit de suivi \niveauexpert}]
Six mois après le rapport, NYANGA annonce que « toutes les recommandations
sont mises en œuvre ». Construisez le plan de l'audit de suivi~: exigences
retestées, preuves attendues, échantillons, et la manière de présenter
l'évolution de $\mu$, $U$ et $H$ par domaine.
\end{boxexo}

\bigskip

\begin{boxinfo}
\textbf{Pour aller plus loin.}
\begin{itemize}
    \item L'ISO~19011~\cite{iso19011_2018}, pour le rapport d'audit et le
          suivi.
    \item T.~E.~Minka Mi Nguidjoï~\cite{minka2026nek}, pour la lecture des
          trois mesures.
\end{itemize}
\end{boxinfo}

\begin{boxinfo}
\textbf{Transition.} Un traitement de NOUN PAY n'a pas encore été audité~:
la notation automatique qui décide des crédits. Profilage, décisions
automatisées, biométrie~: ce sont les défis de la Partie~V, et l'objet du
Chapitre~\ref{chap:profilage}.
\end{boxinfo}
